% Shared preamble for the paper and its supplement (pdflatex).
\PassOptionsToPackage{hyphens}{url}
\usepackage[T1]{fontenc}
\usepackage[utf8]{inputenc}
\usepackage{lmodern}
\usepackage[letterpaper,margin=1in]{geometry}
\usepackage{setspace}
\setstretch{1.12}
\usepackage{amsmath,amssymb,amsthm,mathtools}
\usepackage{booktabs,array,tabularx,longtable}
\usepackage{enumitem}
\usepackage{xcolor}
\usepackage{microtype}
\usepackage[authoryear,round]{natbib}
\usepackage{xr-hyper}
\usepackage[hidelinks]{hyperref}
\urlstyle{same}
\setlength{\emergencystretch}{3em}
\setlist{itemsep=2pt,topsep=4pt}
\newcolumntype{L}[1]{>{\raggedright\arraybackslash\hspace{0pt}}p{#1}}
\newcolumntype{Y}{>{\raggedright\arraybackslash}X}

% ---------------- notation ----------------
\newcommand{\Bset}{\mathfrak B}
\newcommand{\Kspace}{\mathcal K}
\newcommand{\ThI}{\Theta_I}
\newcommand{\ThX}{\Theta_X}
\newcommand{\ThO}{\Theta_O}
\newcommand{\E}{\mathbb E}
\newcommand{\R}{\mathbb R}
\newcommand{\conv}{\operatorname{conv}}
\newcommand{\supp}{\operatorname{supp}}
\newcommand{\vertx}{\operatorname{vert}}
\newcommand{\affdim}{\operatorname{affdim}}
\newcommand{\Var}{\operatorname{Var}}
\newcommand{\Cov}{\operatorname{Cov}}
\newcommand{\sym}{\operatorname{sym}}
\newcommand{\skw}{\operatorname{skew}}
\newcommand{\diag}{\operatorname{diag}}
\newcommand{\hull}{\operatorname{hull}}
\newcommand{\sgn}{\operatorname{sign}}
\newcommand{\one}{\mathbf 1}
\newcommand{\Bt}{\tilde B}
\newcommand{\DD}{\mathrm{DD}}
\newcommand{\CI}{\mathrm{CI}}
\newcommand{\ER}{\mathrm{ER}}
\newcommand{\NM}{\mathrm{NM}}
\newcommand{\pkg}[1]{\textsf{#1}}
\newcommand{\code}[1]{\texttt{#1}}

% ---------------- result environments ----------------
\theoremstyle{plain}
\newtheorem{theorem}{Theorem}
\newtheorem{proposition}{Proposition}
\newtheorem{lemma}{Lemma}
\newtheorem{corollary}{Corollary}
\newtheorem{conjecture}{Conjecture}
\theoremstyle{definition}
\newtheorem{assumption}{Assumption}
\newtheorem{definition}{Definition}
\newtheorem{procedure}{Procedure}
\theoremstyle{remark}
\newtheorem{remark}{Remark}
\newtheorem{example}{Example}

% plain-language reading and sketch of the argument
\newenvironment{reading}{\par\smallskip\noindent\textit{Reading.}\ }{\par\smallskip}
\newenvironment{sketch}{\par\smallskip\noindent\textit{Sketch of the argument.}\ }{\par\smallskip}
% status tag printed after a result's title
\newcommand{\status}[1]{\textup{[\textsf{#1}]}}
